% =========================================================
% glosario.tex
% Definición y generación del glosario y siglas del TFM
% =========================================================
%
% Este fichero contiene:
%   - Definiciones de términos (glosario)
%   - Definiciones de siglas (acrónimos)
%   - Estructura de impresión del glosario en el documento
%
% Se utiliza junto con el paquete:
%   \usepackage[acronym]{glossaries}
%
% ---------------------------------------------------------
% 📌 ¿Cómo usar este fichero?
% ---------------------------------------------------------
%
% 1. Define nuevos términos con:
%
%    \newglossaryentry{clave}{
%        name={Nombre visible},
%        description={Definición del término}
%    }
%
% 2. Define nuevas siglas con:
%
%    \newacronym{clave}{SIGLA}{Nombre completo}
%
% 3. Usa los términos en el documento:
%
%    \gls{clave}        → muestra el término
%
% 4. Usa las siglas:
%
%    \acrshort{clave}   → forma corta (API)
%    \acrlong{clave}    → forma larga (Interfaz de Programación de Aplicaciones)
%
% 5. Genera el glosario automáticamente con:
%
%    \crearglosario
%
% ---------------------------------------------------------
% 📖 Tipos de entradas
% ---------------------------------------------------------
%
% 🔤 Términos (glosario)
%
% Ejemplo:
%   \newglossaryentry{backend}{
%       name={Backend},
%       description={Sistema de servidores y lógica}
%   }
%
% 🔠 Siglas (acrónimos)
%
% Ejemplo:
%   \newacronym{api}{API}{Interfaz de Programación de Aplicaciones}
%
% También puedes definir plural:
%
%   \newacronym[
%     plural=SIGLAS,
%     longplural={Nombre completo en plural}
%   ]{clave}{SIGLA}{Nombre completo}
%
% ---------------------------------------------------------
% ⚙️ Impresión del glosario
% ---------------------------------------------------------
%
% Este fichero incluye dos secciones:
%
% 1. Glosario de siglas
%    - Muestra únicamente acrónimos
%
% 2. Definiciones
%    - Muestra términos definidos con \newglossaryentry
%
% Ambas secciones se insertan mediante:
%
%   % =========================================================
% glosario.tex
% Definición y generación del glosario y siglas del TFM
% =========================================================
%
% Este fichero contiene:
%   - Definiciones de términos (glosario)
%   - Definiciones de siglas (acrónimos)
%   - Estructura de impresión del glosario en el documento
%
% Se utiliza junto con el paquete:
%   \usepackage[acronym]{glossaries}
%
% ---------------------------------------------------------
% 📌 ¿Cómo usar este fichero?
% ---------------------------------------------------------
%
% 1. Define nuevos términos con:
%
%    \newglossaryentry{clave}{
%        name={Nombre visible},
%        description={Definición del término}
%    }
%
% 2. Define nuevas siglas con:
%
%    \newacronym{clave}{SIGLA}{Nombre completo}
%
% 3. Usa los términos en el documento:
%
%    \gls{clave}        → muestra el término
%
% 4. Usa las siglas:
%
%    \acrshort{clave}   → forma corta (API)
%    \acrlong{clave}    → forma larga (Interfaz de Programación de Aplicaciones)
%
% 5. Genera el glosario automáticamente con:
%
%    \crearglosario
%
% ---------------------------------------------------------
% 📖 Tipos de entradas
% ---------------------------------------------------------
%
% 🔤 Términos (glosario)
%
% Ejemplo:
%   \newglossaryentry{backend}{
%       name={Backend},
%       description={Sistema de servidores y lógica}
%   }
%
% 🔠 Siglas (acrónimos)
%
% Ejemplo:
%   \newacronym{api}{API}{Interfaz de Programación de Aplicaciones}
%
% También puedes definir plural:
%
%   \newacronym[
%     plural=SIGLAS,
%     longplural={Nombre completo en plural}
%   ]{clave}{SIGLA}{Nombre completo}
%
% ---------------------------------------------------------
% ⚙️ Impresión del glosario
% ---------------------------------------------------------
%
% Este fichero incluye dos secciones:
%
% 1. Glosario de siglas
%    - Muestra únicamente acrónimos
%
% 2. Definiciones
%    - Muestra términos definidos con \newglossaryentry
%
% Ambas secciones se insertan mediante:
%
%   % =========================================================
% glosario.tex
% Definición y generación del glosario y siglas del TFM
% =========================================================
%
% Este fichero contiene:
%   - Definiciones de términos (glosario)
%   - Definiciones de siglas (acrónimos)
%   - Estructura de impresión del glosario en el documento
%
% Se utiliza junto con el paquete:
%   \usepackage[acronym]{glossaries}
%
% ---------------------------------------------------------
% 📌 ¿Cómo usar este fichero?
% ---------------------------------------------------------
%
% 1. Define nuevos términos con:
%
%    \newglossaryentry{clave}{
%        name={Nombre visible},
%        description={Definición del término}
%    }
%
% 2. Define nuevas siglas con:
%
%    \newacronym{clave}{SIGLA}{Nombre completo}
%
% 3. Usa los términos en el documento:
%
%    \gls{clave}        → muestra el término
%
% 4. Usa las siglas:
%
%    \acrshort{clave}   → forma corta (API)
%    \acrlong{clave}    → forma larga (Interfaz de Programación de Aplicaciones)
%
% 5. Genera el glosario automáticamente con:
%
%    \crearglosario
%
% ---------------------------------------------------------
% 📖 Tipos de entradas
% ---------------------------------------------------------
%
% 🔤 Términos (glosario)
%
% Ejemplo:
%   \newglossaryentry{backend}{
%       name={Backend},
%       description={Sistema de servidores y lógica}
%   }
%
% 🔠 Siglas (acrónimos)
%
% Ejemplo:
%   \newacronym{api}{API}{Interfaz de Programación de Aplicaciones}
%
% También puedes definir plural:
%
%   \newacronym[
%     plural=SIGLAS,
%     longplural={Nombre completo en plural}
%   ]{clave}{SIGLA}{Nombre completo}
%
% ---------------------------------------------------------
% ⚙️ Impresión del glosario
% ---------------------------------------------------------
%
% Este fichero incluye dos secciones:
%
% 1. Glosario de siglas
%    - Muestra únicamente acrónimos
%
% 2. Definiciones
%    - Muestra términos definidos con \newglossaryentry
%
% Ambas secciones se insertan mediante:
%
%   % =========================================================
% glosario.tex
% Definición y generación del glosario y siglas del TFM
% =========================================================
%
% Este fichero contiene:
%   - Definiciones de términos (glosario)
%   - Definiciones de siglas (acrónimos)
%   - Estructura de impresión del glosario en el documento
%
% Se utiliza junto con el paquete:
%   \usepackage[acronym]{glossaries}
%
% ---------------------------------------------------------
% 📌 ¿Cómo usar este fichero?
% ---------------------------------------------------------
%
% 1. Define nuevos términos con:
%
%    \newglossaryentry{clave}{
%        name={Nombre visible},
%        description={Definición del término}
%    }
%
% 2. Define nuevas siglas con:
%
%    \newacronym{clave}{SIGLA}{Nombre completo}
%
% 3. Usa los términos en el documento:
%
%    \gls{clave}        → muestra el término
%
% 4. Usa las siglas:
%
%    \acrshort{clave}   → forma corta (API)
%    \acrlong{clave}    → forma larga (Interfaz de Programación de Aplicaciones)
%
% 5. Genera el glosario automáticamente con:
%
%    \crearglosario
%
% ---------------------------------------------------------
% 📖 Tipos de entradas
% ---------------------------------------------------------
%
% 🔤 Términos (glosario)
%
% Ejemplo:
%   \newglossaryentry{backend}{
%       name={Backend},
%       description={Sistema de servidores y lógica}
%   }
%
% 🔠 Siglas (acrónimos)
%
% Ejemplo:
%   \newacronym{api}{API}{Interfaz de Programación de Aplicaciones}
%
% También puedes definir plural:
%
%   \newacronym[
%     plural=SIGLAS,
%     longplural={Nombre completo en plural}
%   ]{clave}{SIGLA}{Nombre completo}
%
% ---------------------------------------------------------
% ⚙️ Impresión del glosario
% ---------------------------------------------------------
%
% Este fichero incluye dos secciones:
%
% 1. Glosario de siglas
%    - Muestra únicamente acrónimos
%
% 2. Definiciones
%    - Muestra términos definidos con \newglossaryentry
%
% Ambas secciones se insertan mediante:
%
%   \input{Glosario/glosario}
%
% y son llamadas automáticamente desde:
%
%   \crearglosario
%
% ---------------------------------------------------------
% ⚠️ Reglas importantes
% ---------------------------------------------------------
%
% - Cada entrada debe tener una clave única:
%     ❌ backend duplicado
%     ✔ backend, backend_api, backend_sistema
%
% - No usar espacios en las claves:
%     ❌ "base de datos"
%     ✔ base_datos
%
% - Usa nombres descriptivos:
%     ✔ resultset
%     ❌ r1
%
% - Evita repetir definiciones en el texto:
%     usa siempre \gls{}
%
% ---------------------------------------------------------
% 💡 Notas
% ---------------------------------------------------------
%
% - El comando \glsaddall se usa para incluir todos los términos
%   aunque no aparezcan en el texto
%
% - Puedes eliminarlo si solo quieres mostrar términos usados
%
% - El formato visual está controlado por la plantilla
%
% =========================================================

% Ejemplos de entradas de glosario y siglas. Reemplaza o añade las tuyas propias.




% =========================
% Glosario de términos
% =========================

\newglossaryentry{redteam}{
    name={Red Team},
    description={Disciplina de seguridad ofensiva orientada a simular ataques reales para evaluar la capacidad de detección y respuesta de una organización}
}

\newglossaryentry{pentesting}{
    name={Pentesting},
    description={Proceso autorizado de evaluación de la seguridad de un sistema mediante técnicas de ataque controladas}
}

\newglossaryentry{pivoting}{
    name={Pivoting},
    description={Técnica que permite utilizar un sistema comprometido como punto de acceso para alcanzar otros sistemas de una red interna}
}

\newglossaryentry{lateralmovement}{
    name={Movimiento lateral},
    description={Conjunto de técnicas utilizadas por un atacante para desplazarse entre distintos sistemas de una red tras obtener acceso inicial}
}

\newglossaryentry{hostonly}{
    name={Red Host-Only},
    description={Tipo de red virtual aislada que permite la comunicación entre máquinas virtuales y el equipo anfitrión sin exponerlas a redes externas}
}

\newglossaryentry{snapshot}{
    name={Snapshot},
    description={Captura instantánea del estado completo de una máquina virtual que permite restaurarla posteriormente}
}


% =========================
% Glosario de siglas
% =========================

\newacronym{owasp}{OWASP}{Open Worldwide Application Security Project}

\newacronym{ad}{AD}{Active Directory}

\newacronym{vm}{VM}{Máquina Virtual}

\newacronym{dns}{DNS}{Domain Name System}

\newacronym{smb}{SMB}{Server Message Block}

\newacronym{winrm}{WinRM}{Windows Remote Management}

\newacronym{http}{HTTP}{HyperText Transfer Protocol}

\newacronym{https}{HTTPS}{HyperText Transfer Protocol Secure}

\newacronym{php}{PHP}{PHP Hypertext Preprocessor}

\newacronym{sql}{SQL}{Structured Query Language}

\newacronym{mysql}{MySQL}{My Structured Query Language}

\newacronym{apache}{Apache}{Apache HTTP Server}

\newacronym{ssh}{SSH}{Secure Shell}

\newacronym{nat}{NAT}{Network Address Translation}

\newacronym{lan}{LAN}{Local Area Network}

\newacronym{ip}{IP}{Internet Protocol}

\newacronym{cpu}{CPU}{Central Processing Unit}

\newacronym{ram}{RAM}{Random Access Memory}

\newacronym{xfce}{XFCE}{XForms Common Environment}

\newacronym{mitre}{MITRE ATT\&CK}{Adversarial Tactics, Techniques and Common Knowledge}

\newacronym{sqli}{SQLi}{SQL Injection}

\newacronym{xss}{XSS}{Cross-Site Scripting}

\newacronym{dvwa}{DVWA}{Damn Vulnerable Web Application}

\newacronym{cve}{CVE}{Common Vulnerabilities and Exposures}

\newacronym{cwe}{CWE}{Common Weakness Enumeration}

\newacronym{git}{Git}{Sistema de control de versiones distribuido}

\newacronym{github}{GitHub}{Plataforma de alojamiento y gestión de repositorios Git}

\newacronym{api}{API}{Interfaz de Programación de Aplicaciones}


% =========================
% Impresión del glosario
% =========================

\glsaddall

\section*{Glosario de siglas}
\label{sec:glosario-siglas}
\vspace{-1.3cm}

\printglossary[type=\acronymtype, title={}, toctitle={}]

\section*{Definiciones}
\label{sec:glosario-terminos}
\vspace{-1.3cm}

\printglossary[title={}, toctitle={}]

%
% y son llamadas automáticamente desde:
%
%   \crearglosario
%
% ---------------------------------------------------------
% ⚠️ Reglas importantes
% ---------------------------------------------------------
%
% - Cada entrada debe tener una clave única:
%     ❌ backend duplicado
%     ✔ backend, backend_api, backend_sistema
%
% - No usar espacios en las claves:
%     ❌ "base de datos"
%     ✔ base_datos
%
% - Usa nombres descriptivos:
%     ✔ resultset
%     ❌ r1
%
% - Evita repetir definiciones en el texto:
%     usa siempre \gls{}
%
% ---------------------------------------------------------
% 💡 Notas
% ---------------------------------------------------------
%
% - El comando \glsaddall se usa para incluir todos los términos
%   aunque no aparezcan en el texto
%
% - Puedes eliminarlo si solo quieres mostrar términos usados
%
% - El formato visual está controlado por la plantilla
%
% =========================================================

% Ejemplos de entradas de glosario y siglas. Reemplaza o añade las tuyas propias.




% =========================
% Glosario de términos
% =========================

\newglossaryentry{redteam}{
    name={Red Team},
    description={Disciplina de seguridad ofensiva orientada a simular ataques reales para evaluar la capacidad de detección y respuesta de una organización}
}

\newglossaryentry{pentesting}{
    name={Pentesting},
    description={Proceso autorizado de evaluación de la seguridad de un sistema mediante técnicas de ataque controladas}
}

\newglossaryentry{pivoting}{
    name={Pivoting},
    description={Técnica que permite utilizar un sistema comprometido como punto de acceso para alcanzar otros sistemas de una red interna}
}

\newglossaryentry{lateralmovement}{
    name={Movimiento lateral},
    description={Conjunto de técnicas utilizadas por un atacante para desplazarse entre distintos sistemas de una red tras obtener acceso inicial}
}

\newglossaryentry{hostonly}{
    name={Red Host-Only},
    description={Tipo de red virtual aislada que permite la comunicación entre máquinas virtuales y el equipo anfitrión sin exponerlas a redes externas}
}

\newglossaryentry{snapshot}{
    name={Snapshot},
    description={Captura instantánea del estado completo de una máquina virtual que permite restaurarla posteriormente}
}


% =========================
% Glosario de siglas
% =========================

\newacronym{owasp}{OWASP}{Open Worldwide Application Security Project}

\newacronym{ad}{AD}{Active Directory}

\newacronym{vm}{VM}{Máquina Virtual}

\newacronym{dns}{DNS}{Domain Name System}

\newacronym{smb}{SMB}{Server Message Block}

\newacronym{winrm}{WinRM}{Windows Remote Management}

\newacronym{http}{HTTP}{HyperText Transfer Protocol}

\newacronym{https}{HTTPS}{HyperText Transfer Protocol Secure}

\newacronym{php}{PHP}{PHP Hypertext Preprocessor}

\newacronym{sql}{SQL}{Structured Query Language}

\newacronym{mysql}{MySQL}{My Structured Query Language}

\newacronym{apache}{Apache}{Apache HTTP Server}

\newacronym{ssh}{SSH}{Secure Shell}

\newacronym{nat}{NAT}{Network Address Translation}

\newacronym{lan}{LAN}{Local Area Network}

\newacronym{ip}{IP}{Internet Protocol}

\newacronym{cpu}{CPU}{Central Processing Unit}

\newacronym{ram}{RAM}{Random Access Memory}

\newacronym{xfce}{XFCE}{XForms Common Environment}

\newacronym{mitre}{MITRE ATT\&CK}{Adversarial Tactics, Techniques and Common Knowledge}

\newacronym{sqli}{SQLi}{SQL Injection}

\newacronym{xss}{XSS}{Cross-Site Scripting}

\newacronym{dvwa}{DVWA}{Damn Vulnerable Web Application}

\newacronym{cve}{CVE}{Common Vulnerabilities and Exposures}

\newacronym{cwe}{CWE}{Common Weakness Enumeration}

\newacronym{git}{Git}{Sistema de control de versiones distribuido}

\newacronym{github}{GitHub}{Plataforma de alojamiento y gestión de repositorios Git}

\newacronym{api}{API}{Interfaz de Programación de Aplicaciones}


% =========================
% Impresión del glosario
% =========================

\glsaddall

\section*{Glosario de siglas}
\label{sec:glosario-siglas}
\vspace{-1.3cm}

\printglossary[type=\acronymtype, title={}, toctitle={}]

\section*{Definiciones}
\label{sec:glosario-terminos}
\vspace{-1.3cm}

\printglossary[title={}, toctitle={}]

%
% y son llamadas automáticamente desde:
%
%   \crearglosario
%
% ---------------------------------------------------------
% ⚠️ Reglas importantes
% ---------------------------------------------------------
%
% - Cada entrada debe tener una clave única:
%     ❌ backend duplicado
%     ✔ backend, backend_api, backend_sistema
%
% - No usar espacios en las claves:
%     ❌ "base de datos"
%     ✔ base_datos
%
% - Usa nombres descriptivos:
%     ✔ resultset
%     ❌ r1
%
% - Evita repetir definiciones en el texto:
%     usa siempre \gls{}
%
% ---------------------------------------------------------
% 💡 Notas
% ---------------------------------------------------------
%
% - El comando \glsaddall se usa para incluir todos los términos
%   aunque no aparezcan en el texto
%
% - Puedes eliminarlo si solo quieres mostrar términos usados
%
% - El formato visual está controlado por la plantilla
%
% =========================================================

% Ejemplos de entradas de glosario y siglas. Reemplaza o añade las tuyas propias.




% =========================
% Glosario de términos
% =========================

\newglossaryentry{redteam}{
    name={Red Team},
    description={Disciplina de seguridad ofensiva orientada a simular ataques reales para evaluar la capacidad de detección y respuesta de una organización}
}

\newglossaryentry{pentesting}{
    name={Pentesting},
    description={Proceso autorizado de evaluación de la seguridad de un sistema mediante técnicas de ataque controladas}
}

\newglossaryentry{pivoting}{
    name={Pivoting},
    description={Técnica que permite utilizar un sistema comprometido como punto de acceso para alcanzar otros sistemas de una red interna}
}

\newglossaryentry{lateralmovement}{
    name={Movimiento lateral},
    description={Conjunto de técnicas utilizadas por un atacante para desplazarse entre distintos sistemas de una red tras obtener acceso inicial}
}

\newglossaryentry{hostonly}{
    name={Red Host-Only},
    description={Tipo de red virtual aislada que permite la comunicación entre máquinas virtuales y el equipo anfitrión sin exponerlas a redes externas}
}

\newglossaryentry{snapshot}{
    name={Snapshot},
    description={Captura instantánea del estado completo de una máquina virtual que permite restaurarla posteriormente}
}


% =========================
% Glosario de siglas
% =========================

\newacronym{owasp}{OWASP}{Open Worldwide Application Security Project}

\newacronym{ad}{AD}{Active Directory}

\newacronym{vm}{VM}{Máquina Virtual}

\newacronym{dns}{DNS}{Domain Name System}

\newacronym{smb}{SMB}{Server Message Block}

\newacronym{winrm}{WinRM}{Windows Remote Management}

\newacronym{http}{HTTP}{HyperText Transfer Protocol}

\newacronym{https}{HTTPS}{HyperText Transfer Protocol Secure}

\newacronym{php}{PHP}{PHP Hypertext Preprocessor}

\newacronym{sql}{SQL}{Structured Query Language}

\newacronym{mysql}{MySQL}{My Structured Query Language}

\newacronym{apache}{Apache}{Apache HTTP Server}

\newacronym{ssh}{SSH}{Secure Shell}

\newacronym{nat}{NAT}{Network Address Translation}

\newacronym{lan}{LAN}{Local Area Network}

\newacronym{ip}{IP}{Internet Protocol}

\newacronym{cpu}{CPU}{Central Processing Unit}

\newacronym{ram}{RAM}{Random Access Memory}

\newacronym{xfce}{XFCE}{XForms Common Environment}

\newacronym{mitre}{MITRE ATT\&CK}{Adversarial Tactics, Techniques and Common Knowledge}

\newacronym{sqli}{SQLi}{SQL Injection}

\newacronym{xss}{XSS}{Cross-Site Scripting}

\newacronym{dvwa}{DVWA}{Damn Vulnerable Web Application}

\newacronym{cve}{CVE}{Common Vulnerabilities and Exposures}

\newacronym{cwe}{CWE}{Common Weakness Enumeration}

\newacronym{git}{Git}{Sistema de control de versiones distribuido}

\newacronym{github}{GitHub}{Plataforma de alojamiento y gestión de repositorios Git}

\newacronym{api}{API}{Interfaz de Programación de Aplicaciones}


% =========================
% Impresión del glosario
% =========================

\glsaddall

\section*{Glosario de siglas}
\label{sec:glosario-siglas}
\vspace{-1.3cm}

\printglossary[type=\acronymtype, title={}, toctitle={}]

\section*{Definiciones}
\label{sec:glosario-terminos}
\vspace{-1.3cm}

\printglossary[title={}, toctitle={}]

%
% y son llamadas automáticamente desde:
%
%   \crearglosario
%
% ---------------------------------------------------------
% ⚠️ Reglas importantes
% ---------------------------------------------------------
%
% - Cada entrada debe tener una clave única:
%     ❌ backend duplicado
%     ✔ backend, backend_api, backend_sistema
%
% - No usar espacios en las claves:
%     ❌ "base de datos"
%     ✔ base_datos
%
% - Usa nombres descriptivos:
%     ✔ resultset
%     ❌ r1
%
% - Evita repetir definiciones en el texto:
%     usa siempre \gls{}
%
% ---------------------------------------------------------
% 💡 Notas
% ---------------------------------------------------------
%
% - El comando \glsaddall se usa para incluir todos los términos
%   aunque no aparezcan en el texto
%
% - Puedes eliminarlo si solo quieres mostrar términos usados
%
% - El formato visual está controlado por la plantilla
%
% =========================================================

% Ejemplos de entradas de glosario y siglas. Reemplaza o añade las tuyas propias.




% =========================
% Glosario de términos
% =========================

\newglossaryentry{redteam}{
    name={Red Team},
    description={Disciplina de seguridad ofensiva orientada a simular ataques reales para evaluar la capacidad de detección y respuesta de una organización}
}

\newglossaryentry{pentesting}{
    name={Pentesting},
    description={Proceso autorizado de evaluación de la seguridad de un sistema mediante técnicas de ataque controladas}
}

\newglossaryentry{pivoting}{
    name={Pivoting},
    description={Técnica que permite utilizar un sistema comprometido como punto de acceso para alcanzar otros sistemas de una red interna}
}

\newglossaryentry{lateralmovement}{
    name={Movimiento lateral},
    description={Conjunto de técnicas utilizadas por un atacante para desplazarse entre distintos sistemas de una red tras obtener acceso inicial}
}

\newglossaryentry{hostonly}{
    name={Red Host-Only},
    description={Tipo de red virtual aislada que permite la comunicación entre máquinas virtuales y el equipo anfitrión sin exponerlas a redes externas}
}

\newglossaryentry{snapshot}{
    name={Snapshot},
    description={Captura instantánea del estado completo de una máquina virtual que permite restaurarla posteriormente}
}


% =========================
% Glosario de siglas
% =========================

\newacronym{owasp}{OWASP}{Open Worldwide Application Security Project}

\newacronym{ad}{AD}{Active Directory}

\newacronym{vm}{VM}{Máquina Virtual}

\newacronym{dns}{DNS}{Domain Name System}

\newacronym{smb}{SMB}{Server Message Block}

\newacronym{winrm}{WinRM}{Windows Remote Management}

\newacronym{http}{HTTP}{HyperText Transfer Protocol}

\newacronym{https}{HTTPS}{HyperText Transfer Protocol Secure}

\newacronym{php}{PHP}{PHP Hypertext Preprocessor}

\newacronym{sql}{SQL}{Structured Query Language}

\newacronym{mysql}{MySQL}{My Structured Query Language}

\newacronym{apache}{Apache}{Apache HTTP Server}

\newacronym{ssh}{SSH}{Secure Shell}

\newacronym{nat}{NAT}{Network Address Translation}

\newacronym{lan}{LAN}{Local Area Network}

\newacronym{ip}{IP}{Internet Protocol}

\newacronym{cpu}{CPU}{Central Processing Unit}

\newacronym{ram}{RAM}{Random Access Memory}

\newacronym{xfce}{XFCE}{XForms Common Environment}

\newacronym{mitre}{MITRE ATT\&CK}{Adversarial Tactics, Techniques and Common Knowledge}

\newacronym{sqli}{SQLi}{SQL Injection}

\newacronym{xss}{XSS}{Cross-Site Scripting}

\newacronym{dvwa}{DVWA}{Damn Vulnerable Web Application}

\newacronym{cve}{CVE}{Common Vulnerabilities and Exposures}

\newacronym{cwe}{CWE}{Common Weakness Enumeration}

\newacronym{git}{Git}{Sistema de control de versiones distribuido}

\newacronym{github}{GitHub}{Plataforma de alojamiento y gestión de repositorios Git}

\newacronym{api}{API}{Interfaz de Programación de Aplicaciones}


% =========================
% Impresión del glosario
% =========================

\glsaddall

\section*{Glosario de siglas}
\label{sec:glosario-siglas}
\vspace{-1.3cm}

\printglossary[type=\acronymtype, title={}, toctitle={}]

\section*{Definiciones}
\label{sec:glosario-terminos}
\vspace{-1.3cm}

\printglossary[title={}, toctitle={}]
